\section{Selecting the physical coefficient}\label{sec:physical-amplitude}

We prove that the coefficient carried by the real counting roots is the coefficient selected by the two vacuum caps. First we identify the roots at the vacuum twist \(K=i\). At generic nearby parameters, the common-twist trace then isolates their ordinary \(W/J\) term. A separate estimate for the two ends of the cylinder makes the cap coefficient analytic at fixed \(N\); this permits normalization at \(K=i\) and continuation to the physical twist \(K=1\). The right cap keeps twist \(i\), and the left cap keeps twist \(-i\), throughout.

The polynomial construction below works for every \(N\ge1\). The identification with the counting branch and the final coefficient formula use \(N\ge N_0\) from Proposition~\ref{prop:real-counting}.

\subsection{The positive roots at the vacuum endpoint}

Set \(l=\pi/8\) and \(q=e^{2il}\). For this subsection take additive sites \(b_a=i\eta_a\), so their multiplicative coordinates
\(\xi_a=e^{2ib_a}=e^{-2\eta_a}\) are positive. Repeated sites are allowed. For \(z=e^{2iw}\), put
\[
 B(z)=\prod_{a=1}^N\prod_{j=3,4,5}(z-q^j\xi_a).
\]
A polynomial is said to have the missing modes if its coefficients of \(z^{4m+1}\) all vanish.

\begin{proposition}[Vacuum endpoint]\label{prop:positive-vacuum-roots}
For every \(N\ge1\) and every multiset \(\xi_1,\ldots,\xi_N>0\), there is a unique monic real polynomial \(P_0\) of degree \(N\) such that \(H=BP_0\) has the missing modes. The zeros of \(P_0\) are positive and simple.

Write these zeros as \(e^{-2y_j}\), with \(y_1<\cdots<y_N\), set \(t_j=iy_j\), and put \(H_j(z)=H(q^jz)\). For \(K\in\mathbb C^*\), the normalized scalar TQ expression for these roots is the rational function
\begin{equation}\label{eq:vacuum-TQ-quotient}
 F_K(z):=\frac{L(w)}{d_{\rm full}(w)}
 =\frac{K H_1H_6+H_5H_0+K^{-1}H_4H_7}{H_7H_6}.
\end{equation}
It satisfies \(F_i\equiv1\) and \(D_j=-i\) for all \(j\). More precisely, with \(p,r\) as in Section~\ref{sec:counting}, define
\[
 T_{\boldsymbol\xi}(y)
 =\sum_{a=1}^N\int_{-\infty}^{y-\eta_a}p(v)\,dv
   +\sum_{k=1}^N\int_{-\infty}^{y-y_k}r(v)\,dv .
\]
The second sum includes the self term. Then
\[
 T_{\boldsymbol\xi}(y_j)=j-\tfrac12+\tfrac14,\qquad 1\le j\le N .
\]
\end{proposition}

\begin{proof}
We first prove the statement about zeros, including repeated sites. The polynomial \(B\) is real and positive on the positive ray: each factor belonging to one site is
\[
 (x+\xi_a)\bigl((x+\xi_a/\sqrt2)^2+\xi_a^2/2\bigr)>0
 \quad (x>0).
\]
The missing modes impose \(N\) real homogeneous equations on the \(N+1\) coefficients of a real polynomial of degree at most \(N\). Choose a nonzero solution \(P_0\).

The positive-kernel argument belongs to Schoenberg's P\'olya-frequency and variation-diminishing theory~\cite{Schoenberg1950PFII}. We include the transform factorization and the strict confluent-minor argument because repeated sites require derivatives of the kernel. For a site value \(c=\xi_a^{4/3}\), integrate
\[
 H(z)(-z)^{1/3}
 \frac{dz}{z\prod_a(z^4+\xi_a^4)\bigl(c-(-z)^{4/3}\bigr)}
\]
around a keyhole about the positive ray, using the principal logarithm of \(-z\). If the site has multiplicity \(m\), also differentiate in \(c\) through order \(m-1\). The end circles vanish. At \(-\xi_a,q^3\xi_a,q^5\xi_a\) the poles are removable. At \(q\xi_a,q^7\xi_a\) the residues cancel to the full site multiplicity: the missing modes give
\[
 iH_j+H_{j+2}-iH_{j+4}-H_{j+6}=0,
\]
and hence \(H_7=iH_1\) modulo that multiplicity near \(\xi_a\). The local coordinates \((-z)^{4/3}\) at the two poles agree as germs, while the one-third powers give the cancelling phases.

Writing \(r_0=x^{4/3}\), the jump therefore gives orthogonality of \(P_0(x)\), for the positive measure
\[
 d\mu(x)=\frac{B(x)x^{1/3}}{x\prod_a(x^4+\xi_a^4)}\,dx ,
\]
against all the indicated \(c\)-derivatives of
\[
 \frac{c}{c^2+c r_0+r_0^2}.
\]
This family, with any initial collection of derivatives at each repeated parameter, has nonsingular interpolation matrices at distinct positive arguments. To see strictness, pass to logarithmic coordinates and remove positive factors. The kernel becomes \((1+2\cosh(u-v))^{-1}\). Its bilateral Laplace transform is proportional to
\[
 \frac{\sin(\pi t/3)}{\sin(\pi t)}
 \quad\text{and hence to}\quad
 \prod_{k=1,2}\Gamma((k+t)/3)\Gamma((k-t)/3),
\]
by the substitution to a positive half-line, the partial-fraction Beta integral, and the reflection formula. Thus the kernel is a positive constant times the convolution of the four kernels
\[
 f_{k,\pm}(s)=\exp\bigl(\pm ks-\exp(\pm3s)\bigr),\qquad k=1,2.
\]
Each matrix \(f_{k,\pm}(u-v)\) has strictly positive ordered minors, including confluent column minors with derivative orders starting at zero. After stripping positive factors, it is \(\exp(-XY)\) with \(X,Y\) oppositely ordered. The distinct row exponentials form an extended Chebyshev system in the column variable, by Rolle induction after division by the first exponential. Confluence gives the Vandermonde sign, which is positive in these orders. Integral determinant multiplication on ordered arguments preserves this strictness under convolution; differentiation under the integrals is valid because the differentiated kernels retain integrable exponential decay.

If \(P_0\) had fewer than \(N\) sign changes on the positive ray, take one more of the preceding initial derivatives than the number of sign changes. The interpolation determinant vanishing at those change points supplies a function whose signs alternate strictly on their complementary intervals. Its product with \(P_0\) has one strict sign, contradicting orthogonality. Thus \(P_0\) has \(N\) positive simple zeros and degree \(N\). Two projectively different solutions would have a nonzero linear combination of degree below \(N\), which is impossible by the same argument. This proves uniqueness after making \(P_0\) monic.

We next identify the vacuum TQ function. Formula~\eqref{eq:vacuum-TQ-quotient} follows term by term from the shifted factors in \(L/d_{\rm full}\); the constants are fixed, for example, at \(z=0\), where the three terms are \(K,1,K^{-1}\). The same quotient formula holds for any degree-\(N\) root polynomial with nonzero constant term. For the constructed \(H\), at every site and to its full multiplicity the eight germs have the form
\[
 (H_0,\ldots,H_7)=(a,b',c',0,0,0,c'+ia,ib').
\]
This follows from the three zeros supplied by \(B\) and the missing-mode relation. Substitution, also after each cyclic shift, shows that
\[
 iH_1H_6+H_5H_0-iH_4H_7-H_7H_6
\]
vanishes at all eight rotations of every site, to their multiplicities. These are \(8N\) nonzero zeros counted with multiplicity. The polynomial has degree at most \(8N\) and also vanishes at zero, so it is identically zero. Hence \(F_i=1\).

The poles at the shifted roots in \(L\) must consequently cancel. Their first pole coefficients are nonzero units times \(i+D_j\), as in \eqref{eq:common-twist-pole-coefficients}. The units are nonzero because the roots and sites are positive and the roots are distinct; an unshifted equality between a root and a site still gives the nonzero factors \([l]\) and \([2l]\). Thus \(D_j=-i\).

It remains to determine the integer labels of these phases. The integrated masses of \(p\) and \(r\) are \(3/4\) and \(1/4\). The product defining \(D_j\), without a self factor, starts at \(q\) as \(y\to-\infty\), while the self primitive at zero is \(1/8\). The phase derivative formulas in Section~\ref{sec:counting} therefore give
\[
 D_j=\exp\bigl(2\pi iT_{\boldsymbol\xi}(y_j)\bigr).
\]
We determine the labels first when the site scales are very widely separated.

Let one site be \(\epsilon\to0\) while the others remain fixed. Normalize the coefficients of \(H\) to be bounded with a nonzero projective limit. That limit is divisible by \(z^3B_{\rm red}\), where \(B_{\rm red}\) omits this site. Write \(H=\sum h_mz^m\). Evaluation at \(q^{3,4,5}\epsilon\), using \(h_1=0\), shows
\[
 \max\{|h_0|,\epsilon^2|h_2|,\epsilon^3|h_3|\}=O(\epsilon^4).
\]
Indeed the matrix of \(1,z^2,z^3\) at \(q^{3,4,5}\) is invertible; its relevant symmetric sum is \(1+\sqrt2\). Hence the limit is divisible by \(z^4\). After dividing by \(z^4\), the missing modes and the uniqueness already proved for \(N-1\) sites identify the limit as \(H_{\rm red}\), up to scale; for the empty reduced list when \(N=1\), set \(H_{\rm red}=1\), and the assertion is immediate. It has full degree. Consequently exactly one positive zero of \(P_0\) tends to zero, and the others tend to the reduced zeros. Since \([z]H=0\), the sum of the inverses of all zeros of \(H\) is zero; the three new site zeros contribute \(-(1+\sqrt2)/\epsilon\). The small root is therefore asymptotic to \(\epsilon/(1+\sqrt2)\).

Iterating this separation gives configurations with one root comparable to each site and arbitrarily large gaps between their scales. In increasing rapidity order, there
\[
 T_{\boldsymbol\xi}(y_j)\in
 j-1+[1/8,7/8]+o(1).
\]
Since \(D_j=-i\), the only possible phase label in this interval is \(j-1+3/4\). Finally, the monic solution varies continuously in the positive sites by projective uniqueness and full degree. Its ordered zeros remain simple and positive, so the integer phase labels cannot change. This proves the stated counting equations for all positive sites.
\end{proof}

At homogeneous sites \(\xi_a=1\), these phase labels form a contiguous ordered solution at \(\delta=1/4\). The containment estimate in the proof of Proposition~\ref{prop:real-counting} puts every such solution in its quantile rectangle, where uniqueness identifies it with that proposition's solution. For \(N\ge N_0\), the implicit function theorem then continues the same roots along
\[
 K=e^{i\theta},\qquad \delta=\theta/(2\pi),\qquad
 0\le\theta\le\pi/2.
\]
Indeed \(T(y_j)=j-\tfrac12+\delta\) is equivalent to \(D_j=-K\). The roots remain distinct and positive in multiplicative coordinates. The branch is analytic locally in \(K\) and under small complex changes of the additive sites.

\subsection{The isolated ordinary term at generic parameters}

Return to additive site notation \(\mathbf b\). On the selected local root branch write
\[
 a(\mathbf b,K)=\frac{W}{J},\qquad
 J=\det\bigl(\partial_{t_j}\log D_i\bigr),
\]
and let \(F(u;\mathbf b,K)\) be its normalized TQ function. On a solution it can be evaluated as \(P(u)/d_{\rm full}(u)\), using the Laurent polynomial part from Section~\ref{sec:common-twist}; this expression has already cancelled the apparent poles at shifted roots.

\begin{lemma}[Isolation of the ordinary trace term]\label{lem:ordinary-physical-term}
Fix \(N\ge N_0\). In a sufficiently small neighborhood of the homogeneous vacuum parameters, take generic sites, a generic common twist \(K\) for which the toric ends in Proposition~\ref{prop:common-twist} are excluded, and a generic row argument \(u\) near \(l\). For all sufficiently large \(h\), the finite exponential-polynomial expansion of
\[
 G_h=\Lambda\bigl(T^{K^{-1}}(u;\mathbf b)\bigr)^h\Psi
\]
contains the term
\[
 \frac{N!a(\mathbf b,K)}{\tau(1)}\,F(u;\mathbf b,K)^h .
\]
Its coefficient is nonzero, and no other term has that exponential base. Here \(\tau(1)\) is nonzero generically and is independent of \(K\).
\end{lemma}

\begin{proof}
Theorem~\ref{thm:row-moments} and Corollary~\ref{cor:common-twist-jets} apply on the stated generic locus. The jet functional is fixed as \(u,h\) vary; only the Laurent polynomial insertion \((P(u)/d_{\rm full}(u))^h\) changes. Differentiating its powers gives a finite sum of polynomials in \(h\) times exponential bases. Terms with zero base disappear for sufficiently large \(h\). At every jet support point the cleared TQ relation gives \(L=P\), so each remaining base is the value at \(u\) of that point's exact rational TQ function. A generic \(u\) separates nonidentical rational functions.

We show that no other root polynomial has the same rational TQ function as the selected one at generic nearby parameters. In \eqref{eq:vacuum-TQ-quotient}, normalize the two corresponding polynomials \(H\) at zero. Multiplication of \(H\) by any invertible series in \(z^4\) leaves the quotient unchanged: every product in its numerator and denominator has one even and one odd shift. Suppose lower coefficients of two such series have been matched. At degree \(m\), equality of the quotients determines their relative coefficient with multiplier
\[
 \mu_m(K)=
 K(a_m-a_m^7)+(a_m^5+1-a_m^7-a_m^6)
       +K^{-1}(a_m^4-a_m^6),\qquad a_m=q^m .
\]
For \(4\nmid m\), its coefficient \(a_m-a_m^7=a_m-a_m^{-1}\) of \(K\) is nonzero. Thus, after excluding finitely many values of \(K\), all these multipliers are nonzero. At degrees divisible by four use a series in \(z^4\) to match the coefficients; at every other degree the triangular equation then matches them. The quotient of the two root polynomials is therefore invariant under \(z\mapsto iz\).

The selected roots near the vacuum remain separated from their three quarter-turns. If that rational quotient had a pole at a selected root, quarter-turn invariance would require poles at all its rotations, where its denominator has no zero. It has no finite poles, and the degree bound makes it constant. The root polynomials are equal.

At generic nearby sites the selected solution lies in the regular configuration locus and is transverse. Its local contribution to the continued trace is the ordinary \(W/J\) evaluation by Corollary~\ref{cor:common-twist-jets}. The \(N!\) labelings have the same evaluation and Jacobian, and the preceding uniqueness excludes every other support polynomial from their exponential base. The factors in \(W\) are nonzero there, and \(J\ne0\) by continuation from Proposition~\ref{prop:real-counting}. Finally, the independent-twist residue formula with the empty insertion makes \(\tau(1)\) independent of the twists and nonzero generically, as used in the distinct-site normalization. This proves the lemma.
\end{proof}

The selected quotient \(a=W/J\) also extends across an unshifted root--site coincidence on the positive slice. Such a point was excluded from the deliberately small regular configuration locus, but \(S(0)=-1\), \(V(0)=1/[l]^2\), and the marked-pole site numerators retain nonzero \([l][2l]\) factors there. With distinct roots and \(J\ne0\), the equations and \(W\) are holomorphic, and the implicit function theorem gives the same local inverse branch as the limits of generic independent-twist points. Along the whole homogeneous real branch \(W\ne0\) and \(J\ne0\), so \(a\) is analytic and nonzero locally at every point.

The finite-jet representation in Lemma~\ref{lem:ordinary-physical-term} is used only at generic common twists where the toric ends have been excluded. The endpoint identity below will instead concern the analytic continuation of this selected quotient and of the independently defined cap coefficient.

\subsection{The fixed-\(N\) annular gap and the cap ends}

Let \(\mathsf M_N(\mathbf b,u,K)\) be the finite annular diagram row of Section~\ref{sec:cylinder}, with local arguments \(u-b_a\), contractible loop weight zero, and winding weight \(n=K+K^{-1}\). Let \(\mathsf A_N\) be its principal submatrix on nonempty annular states. At homogeneous sites and \(K=e^{i\theta}\), its weights are nonnegative for \(l\le u<2l\) and \(0\le\theta\le\pi/2\).

We describe precisely the finite sums at the ends of \(G_h\). Expand the caps in their finite disk-matching bases and the rows in tile diagrams, summing the orientations of each connected component. A horizontal cut is empty when all its quantum ports are vacant. Let \(E^-_a\) be the sum from the bottom cap through \(a\) rows to an empty cut, with every earlier cut nonempty. Let \(E^+_b\) be the sum from an empty cut through \(b\) rows to the top cap, with every later cut nonempty. Set \(E^-_0=E^+_0=1\), the empty coefficients of the normalized caps. Let \(C_h\) be the sum of diagrams in \(G_h\) with every cut nonempty, including the two end cuts. These are finite sums whenever the cap entries are regular.

Put \(Z_m^{00}=\langle0|\mathsf M_N^m|0\rangle\). At an empty cut the orientation sums split, and the portion between two empty cuts has exactly the ordinary \(0,n\) loop weights. Choosing the first and last empty cuts therefore gives the exact identity
\begin{equation}\label{eq:cap-end-decomposition}
 G_h=C_h+\sum_{\substack{a,b\ge0\\a+b\le h}}
 E^+_b\,Z_{h-a-b}^{00}\,E^-_a .
\end{equation}
The case \(a+b=h\) has a single empty cut and \(Z_0^{00}=1\). At \(h=0\), null-loop contraction gives \(C_0=0\), consistently with \(\Lambda\Psi=1\).

\begin{proposition}[Fixed-\(N\) annular gap and cap estimate]\label{prop:fixed-N-cap-estimate}
Fix \(N\ge N_0\). There is \(\varepsilon_N>0\) with the following properties. For every
\[
 p_0=(\mathbf0,u_0,e^{i\theta_0}),\qquad
 l\le u_0\le l+\varepsilon_N,\quad 0\le\theta_0\le\pi/2,
\]
the matrix \(\mathsf M_N(p_0)\) has a simple dominant eigenvalue \(\lambda(p_0)>0\), and
\[
 \rho\bigl(\mathsf A_N(p_0)\bigr)<\lambda(p_0).
\]
The right cap \(\Psi\) at twist \(i\) and the left cap \(\Lambda\) at twist \(-i\) are holomorphic in the sites near \(\mathbf0\). At homogeneous sites they are limits of rows from vacuum-empty ends at their respective twists, with reversal and spin negation on the upper end as in the definition of \(\Lambda\).

There is a complex neighborhood \(\mathcal U\) of \(p_0\), positive constants \(C,\gamma,\beta\), and holomorphic functions \(\lambda,c_{\rm phys}\) on \(\mathcal U\), with
\[
 \gamma<\beta<\inf_{p\in\mathcal U}|\lambda(p)|,
\]
such that, for \(p=(\mathbf b,u,K)\in\mathcal U\) and all \(h\ge0\),
\begin{equation}\label{eq:fixed-N-cap-bounds}
 \begin{aligned}
 |E^-_h(p)|+|E^+_h(p)|+|C_h(p)|&\le C\gamma^h,\\
 |G_h(p)-c_{\rm phys}(p)\lambda(p)^h|&\le C\beta^h.
 \end{aligned}
\end{equation}
Here \(\lambda(p)\) continues the stated annular eigenvalue. At \(K=i\) in a vacuum neighborhood, \(\lambda=1\) and \(c_{\rm phys}=1\). The gaps, constants and neighborhoods in this proposition may depend on \(N\).
\end{proposition}

\begin{proof}
We first establish the annular gap. When \(n>0\), the nonnegative annular matrix is primitive. From a nonempty state, connect cyclically consecutive occupied ends of different arcs, carrying all other ends straight, until only one arc remains. Two single turns across the intervening empty interval perform each merger. Close the last arc across its complementary gap, producing a winding loop with positive weight \(n\). Conversely, reverse deletions of minimal disk-side arcs to build every state from the empty state, using the two output turns. These moves have positive weight also at \(u=l\), and the empty state has a positive self-step.

For completeness, the finite matrix fact used here is as follows. If a nonnegative matrix \(M\) has a strictly positive power, its spectral radius \(r>0\) has a strictly positive eigenvector: normalize
\(\sum_{h\ge0}t^{-h}M^h{\bf1}\) as \(t\downarrow r\). Its norm diverges because the entry sum of \(M^h\) dominates \(r^h\), and a limit is the required eigenvector. Conjugating by that vector and dividing by \(r\) makes \(M\) row-stochastic. Equality in the sup-norm averaging inequality for a positive power forces any unit-circle eigenvector to be constant, while bounded powers exclude Jordan growth. Thus \(r\) is simple and strictly dominant. Deleting a state strictly decreases the spectral radius, because a sufficiently high power of the restricted substochastic matrix loses mass in every row. Applied to the empty state, this gives \(\rho(\mathsf A_N)<\lambda\).

At \(n=0,u=l\), Section~\ref{sec:cylinder} proves
\[
 (\mathsf M_N(\mathbf0,l,i))^h
 \longrightarrow |\Psi_{\rm ann}\rangle\langle0|.
\]
The empty functional is fixed. In finite dimension this convergence gives the simple eigenvalue \(1\) and spectral radius below \(1\) on the nonempty block. Continuity supplies the same gaps for \(u\) near \(l\) at \(n=0\). Compactness of \(0\le\theta\le\pi/2\) now supplies an \(\varepsilon_N>0\) for the stated interval. This compactness is for the fixed matrix size \(N\).

We next obtain regular caps. At zero loop weights, forget the annular homotopy of every arc and send the resulting disk matching to its tensor in \(D\). This map is onto: every disk matching can be drawn in a cut cylinder. It intertwines the annular row with the spin row of twist \(i\). Indeed the auxiliary--spectator contraction of Section~\ref{sec:vacuum} straightens an arc or closes it with weight zero; the exterior contraction has the same alternatives and phases. Forgetting which side contains the annular hole is the only change in connectivity. The image of the nonempty annular subspace is \(D'\), the subspace of empty coefficient zero. Hence the restriction of the spin row to \(D'\) at homogeneous \(u=l\) is a quotient of the nonempty annular block and has spectral radius below one.

The normalized fixed vector of this test row in \(D\) is therefore unique and holomorphic in nearby sites, by solving its finite linear equations on \(D'\). Commutation and empty normalization make it the simultaneous vacuum, so it agrees with \(\Psi\) at generic sites. Reversal, transposition and spin negation give the same conclusion for \(\Lambda\) at twist \(-i\). Projecting the preceding annular limit proves the asserted limits from the lower empty end; applying the same reversal and negation proves the upper limit.

We now bound the end sums. Work first at a real point \(p_0\) in the proposition, and abbreviate \(\mathsf A=\mathsf A_N(p_0)\). Choose
\[
 \rho(\mathsf A)<a<\lambda(p_0).
\]
There is a constant \(C_A\) with \(\|\mathsf A^m\|\le C_Aa^m\) for all \(m\ge0\). Use any fixed submultiplicative matrix norm dominating entries.

Consider a tile history with nonempty cuts. Complete its bottom occupation by a noncrossing annular arch pattern below that cut, chosen from the bottom cap matching when there is one. Mark the added arches. Propagate both connectivity and the marks upward. A closure using a marked arch consumes at least one mark, so there are at most \(N\) rows with such closures. Every other closed component is wholly in the bulk and has its ordinary weight zero or \(n\). The original occupation remains nonempty at each retained cut, so every ordinary intervening row is counted by the principal block \(\mathsf A\).

Choose a fixed \(M_*\ge\max(1,n)\), and let \(\mathsf X\) be the principal matrix on the same nonempty states obtained from the same local tiles by giving every closed loop weight \(M_*\). It has finite entries, since a row and its state set are finite. After the bounded cap orientation factors are accounted for, it dominates every row with a marked closure, including any other closures in that row. Given the initial completion and a tile history, the marks propagate uniquely; forgetting them introduces only the finite sum over initial cap choices, with no multiplicity at each row. Define the finite matrix
\[
 \mathsf U_m=
 \sum_{r=0}^{\min(N,m)}
 \ \sum_{\substack{m_0,\ldots,m_r\ge0\\m_0+\cdots+m_r=m-r}}
 \mathsf A^{m_r}\mathsf X\mathsf A^{m_{r-1}}
       \cdots\mathsf X\mathsf A^{m_0},
 \qquad m\ge0.
\]
The term with \(r=0\) is \(\mathsf A^m\). This sum records the at most \(N\) exceptional rows and all intervening nonempty-block powers.

At real phases, sum the orientations of each component before taking absolute values. Bulk loops retain their \(0,n\) weights. Components meeting a cap have bounded orientation sums for fixed \(N\), even if they wind many times, because their phases have modulus one and there are at most \(2N\) boundary strands. The finite cap coefficients can be absorbed into a constant. It follows that
\[
 |C_m|\le C_N\|\mathsf U_m\|,\qquad
 |E^-_m|+|E^+_m|\le C_N\|\mathsf U_{m-1}\|
 \quad(m\ge1).
\]
For \(E^-_m\) strip the final nonempty-to-empty row and bound that one transition by a finite constant. For \(E^+_m\) strip its initial empty-to-nonempty row; its actual output arch state supplies a completion for the remaining rows. Marking those arcs is a valid upper bound, including when a later closure is an ordinary bulk loop. The top cap then contributes only the bounded terminal strand factors.

For a fixed \(r\), the norm of each word in \(\mathsf U_m\) is at most
\(C_A^{r+1}\|\mathsf X\|^r a^{m-r}\), and there are \(\binom mr\) choices of the exceptional row positions. Since \(a>0\) and \(r\le N\),
\begin{equation}\label{eq:cap-word-bound}
 \|\mathsf U_m\|\le C_N(m+1)^N a^m .
\end{equation}

The same estimate holds in a sufficiently small complex neighborhood with a slightly larger exponential rate. To see this, majorize the moduli of the local tile weights by nearby positive weights, and use \(|K+K^{-1}|\) for an ordinary winding loop. The resulting nonnegative block \(\widehat{\mathsf A}\) tends to \(\mathsf A\) as the neighborhood shrinks; the corresponding exceptional matrix remains bounded. A fixed resolvent circle gives a uniform power bound for \(\widehat{\mathsf A}\) with rate still below \(|\lambda|\). A boundary component can acquire an additional seam modulus under a complex perturbation. If the size of that perturbation is \(\epsilon\), its total cost through \(m\) rows is at most \(\exp(O(\epsilon Nm))\), since only the finitely many boundary strands are involved. Choose the neighborhood small enough that this cost, and then the polynomial in \eqref{eq:cap-word-bound}, can be absorbed into a rate \(\gamma<\inf|\lambda|\). The caps are bounded there by their holomorphy. This proves the first estimate in \eqref{eq:fixed-N-cap-bounds}.

Finally define the convergent end series
\[
 \mathcal E^\pm(z)=\sum_{m\ge0}E^\pm_m z^m,\qquad
 \mathcal C(z)=\sum_{m\ge0}C_mz^m .
\]
They are holomorphic for \(|z|<1/\gamma\), jointly with the parameters after shrinking the neighborhood. The middle series is
\[
 \mathcal Z^{00}(z)
 =\langle0|(I-z\mathsf M_N)^{-1}|0\rangle
 =\frac{d}{1-z\lambda}+\mathcal R(z),
 \qquad d=\langle0|\Pi_\lambda|0\rangle,
\]
where \(\Pi_\lambda\) is the spectral projector. The annular gap makes \(\mathcal R\) holomorphic in a disk of radius strictly larger than \(1/|\lambda|\), uniformly on a smaller neighborhood. Summing \eqref{eq:cap-end-decomposition} gives
\[
 \sum_{h\ge0}G_hz^h
 =\mathcal C(z)+\mathcal E^+(z)\mathcal Z^{00}(z)\mathcal E^-(z).
\]
Its only possible pole in a suitable larger disk is at \(z=\lambda^{-1}\), of order at most one. The coefficient of \((1-z\lambda)^{-1}\) is
\[
 c_{\rm phys}
 =d\,\mathcal E^+(\lambda^{-1})\mathcal E^-(\lambda^{-1}),
\]
which is holomorphic in the parameters. Subtracting this polar term and applying Cauchy's coefficient bound on a fixed circle beyond \(z=\lambda^{-1}\) gives the second estimate in \eqref{eq:fixed-N-cap-bounds}, with \(\gamma<\beta<\inf|\lambda|\).

At \(K=i\), the bulk spin twist \(K^{-1}=-i\) agrees with the left vacuum twist. Thus \(\Lambda T^{-i}(u)=\Lambda\) and \(G_h=\Lambda\Psi=1\). The empty annular functional is fixed at zero loop weight for all nearby local arguments, so the annular branch has \(\lambda=1\). The estimate forces \(c_{\rm phys}=1\).
\end{proof}

\subsection{Continuation to the homogeneous physical parameters}

The generic trace term and the cap estimate now identify the same analytic coefficient. The following statement retains the distinct bulk and cap twists.

\begin{proposition}[The cap-selected coefficient]\label{prop:physical-coefficient}
Fix \(N\ge N_0\), homogeneous sites, and \(0\le\theta\le\pi/2\). Evaluate \(W,J\) at the real roots of Proposition~\ref{prop:real-counting} with \(\delta=\theta/(2\pi)\), and write the quotient as \((W/J)_\theta\). Put \(F_\theta(u)=F(u;\mathbf0,e^{i\theta})\). In the cap matrix element with bulk spin twist \(K^{-1}=e^{-i\theta}\), right cap twist \(i\), and left cap twist \(-i\), the eigenvalue selected in Proposition~\ref{prop:fixed-N-cap-estimate} is the TQ value \(F_\theta(l)\). Its coefficient is
\begin{equation}\label{eq:physical-coefficient}
 c_{\rm phys}(\theta)=\frac{(W/J)_\theta}{(W/J)_{\pi/2}} .
\end{equation}
In particular \(c_{\rm phys}(\pi/2)=1\).
\end{proposition}

\begin{proof}
Start near the vacuum with a row argument \(l<u<l+\varepsilon_N\). At homogeneous sites and \(K=i\), the selected TQ function is \(1\) by Proposition~\ref{prop:positive-vacuum-roots}, while the cap estimate has \(\lambda=1\) and an error rate below one. Shrink the neighborhood so that \(|F(u;\mathbf b,K)|>\beta\). For generic parameters there, Lemma~\ref{lem:ordinary-physical-term} gives a nonzero isolated term with base \(F(u;\mathbf b,K)\). The generating function of an exponential polynomial has a pole at the reciprocal of each uncancelled nonzero base. The cap generating function has no such pole in this disk except possibly at \(\lambda^{-1}\). Therefore \(F(u;\mathbf b,K)=\lambda(\mathbf b,u,K)\), and comparison of the simple-pole coefficients gives
\[
 c_{\rm phys}(\mathbf b,u,K)
 =\frac{N!a(\mathbf b,K)}{\tau(1)}
\]
at these generic parameters.

Take the limit of this equality to \(K=i\) at generic nearby sites. The cap coefficient tends to \(1\), and \(a(\mathbf b,K)\) tends to the nonzero value \(a(\mathbf b,i)\). Since \(\tau(1)\) is independent of \(K\), the common normalization is thereby identified, giving
\[
 c_{\rm phys}(\mathbf b,u,K)
 =\frac{a(\mathbf b,K)}{a(\mathbf b,i)}
\]
near the vacuum by analytic continuation. This limiting argument uses the generic trace identity only before taking the limit. No representation by jets supported only in the root torus is asserted at \(K=i\) or \(K=1\), which can be toric-end resonances.

Now restrict the sites to the homogeneous value and continue along \(K=e^{i\theta}\) from \(i\) to \(1\). Proposition~\ref{prop:real-counting} supplies locally analytic roots and nonzero \(J\), and the preceding observation gives \(W\ne0\), all along this interval. Proposition~\ref{prop:fixed-N-cap-estimate} supplies locally analytic annular eigenvalues and cap coefficients there; the simple positive eigenvalue selects the same branch on overlaps. The identity theorem therefore continues both \(F=\lambda\) and the coefficient ratio through the interval. Continuity in the row argument includes \(u=l\), where \(P(u)/d_{\rm full}(u)\) is regular because \(d_{\rm full}(l)=([4][3])^N\ne0\). This proves \eqref{eq:physical-coefficient}.
\end{proof}

The next section estimates this quotient uniformly as \(N\to\infty\). The annular gaps and analytic neighborhoods used here were established separately for each fixed \(N\); they supply the exact coefficient identity and make no uniform claim about those gaps.
